\cotefile{157-5}
\scanpage{325}\mspage{241}
% l. 1–5
Soit donc $\unsure{\xi|\zeta}_{\alpha}=(i_{\alpha},u'_{\alpha})$ \inserted{$(\alpha\in A,\operatorname{Card} A\leq\unsure{\pi|\kappa})$} une famille d'éléments de $i'\backslash I$;
\diagram{\begin{tikzcd}[ampersand replacement=\&,column sep=small]
 \& i'_{\alpha}\arrow[rr,bend left=12,"w_{\alpha}"]\arrow[d,no head] \& \deleted{\unsure{f}}j\arrow[r,"\deleted{\unsure{w}}","v"'] \& k\\
 i'\arrow[r,"u'_{\alpha}"]\arrow[rr,bend right=17,"v'"',no head] \& f(i\unsure{'}_{\alpha})\arrow[r,"u''_{\alpha}",no head]\& f(\unsure{j|j'})\arrow[r,"f(v)"]\& f(k)
\end{tikzcd}}
comme $I'$ satisfait à $\mathrm{PF}1_{\pi}$ (cf \unsure{4.9.6.9.7|4.9.6.9.2|4.9.6.9.1}), il existe $j'$ dans $I'$ et des flèches $u''_{\alpha}:f(i'_{\alpha})\to\deleted{\unsure{f(}}j'\deleted{\unsure{)}}$, telles que $u''_{\alpha}u'_{\alpha}=v'$ indépendamment de $\alpha$. \deleted{\unsure{v}} Comme $f$ satisfait
% l. 6–9
$\mathrm{F}1$, OPS que $j'$ est de la forme $f(j)$. \unsure{Appliquer|Appliquons} $\mathrm{F}3^{\bullet}_{\pi}$ aux \deleted{$u''_{\alpha}$ et} $i_{\alpha},j,u''_{\alpha}$, on \unsure{trouve|tr.} un $v:j\longrightarrow k$ et des $w_{\alpha}:\unsure{d_{\alpha}|i_{\alpha}}\unsure{\mathbin{-}|\longrightarrow}k$, tels que
\[
 \underbrace{f(v)u''_{\alpha}}_{\unsure{\text{indép.}|\text{nat}|\text{nât}}}=f(w_{\alpha})\qquad\forall\alpha\in A,
\]
% l. 10–13
d'où
\[
 \deleted{f(v)u''_{\alpha}}\quad f(w_{\alpha})u'_{\alpha}
 = f(v)\underbrace{u''_{\alpha}u'_{\alpha}}_{v'}=f(v)v'\quad\text{indép. des }\alpha,
\]
soit $\unsure{w'}$, donc les $w_{\alpha}$ définissent une flèche des $\xi_{\alpha}$ dans $\xi=(k,w')$, qed.

% l. 14–18
$b''_{\pi}\Longrightarrow\unsure{b}_{\pi}$ \inserted{si $f$ est fidèle}. On est \unsure{ramené} à prouver que si $b''_{\pi}$ est satisfait, \inserted{$f$ fidèle,} alors \deleted{\unsure{on a la condition} $\mathrm{F}3$.} \inserted{les} catégories $i'\backslash I$, dont on sait déjà qu'elles satisfont $\mathrm{PF}1_{\pi}$, satisfont également $\mathrm{PF}2_{\pi}$.

% l. 19–26
\diagram{\begin{tikzcd}[ampersand replacement=\&,column sep=small]
 \& i\deleted{\unsure{\alpha}}\arrow[r,shift left=2,"u_{\alpha}"]\arrow[r]\arrow[r,shift right=2]\arrow[d,no head]\& j\arrow[r,"\unsure{v.}",dotted]\& k\\
 i'\arrow[r,"u'"]\arrow[rr,bend right=15,"v'"',no head]\& f(i\unsure{_{\alpha}|{}})\arrow[r,shift left=2]\arrow[r]\arrow[r,shift right=2]\& f(j)\arrow[r]\& f(k)
\end{tikzcd}}
Soient donc $\xi\unsure{_{1}}=(i,u')$, $\eta=(j,v')$ deux objets de $i'\backslash I$, et $(u_{\alpha})_{\alpha\in A}$ une famille de flèches $\xi\longrightarrow\eta$ \inserted{(avec $\operatorname{Card} A\leq\pi$)}, donc des flèches dans $I$ telles que
\[
 \underline{f(u_{\alpha})u'=v'}\qquad\forall\alpha\unsure{\in}A.
\]
Comme $I'$ est \unsure{grde|grdes|gr des|qv des} devant $\pi$, on peut égaliser par $f(j)\longrightarrow k'$, et comme \deleted{$\unsure{f}$} $k'\backslash I\ne\varnothing$, OPS $k'$ de la forme $f(k)$. Or
% archive sheet number: 324
\scanpage{326}\mspage{242}
% l. 1–3
d'après l'implication $b''\Rightarrow a$\inserted{$\Rightarrow c)$} de la \unsure{partie} $1^\circ)$ (cf page 235) et comme $b''_{\unsure{\pi|\kappa}}\Rightarrow b''$, on sait \deleted{\unsure{alors}} \unsure{déjà} que $f$ satisfait à $\mathrm{F}3$.
% l. 4–7
En l'appliquant à $j,k$ et $f(j)\longrightarrow f(k)$, OPS que $f(j)\longrightarrow f(k)$ provient de $v:j\unsure{\mathbin{-}|\longrightarrow}k$. On aura donc $\deleted{f(\unsure{vu_{\alpha}|vu'_{\alpha}})=f(}\; f(v)f(u_{\alpha})=f(v)f(u_{\beta})$ $\forall\alpha,\beta\in A$, donc $f(vu_{\alpha})=f(vu_{\beta})$.
% l. 8–11
Comme $f$ est fidèle, cela \unsure{indique|implique} que $vu_{\alpha}=vu_{\beta}$, i.e. $v$ égalise les $u_{\alpha}$. D'ailleurs, $v$ peut être considéré comme une flèche \uline{dans} $i'\backslash I$, de $(j,v')$ dans $(k,f(v)v')$. Qed.

% l. 12–15
\uline{4.9.6.10.} Quand $I$ est \uline{ordonné}, alors les quatre conditions $\unsure{a}_{\pi}\unsure{,|\Rightarrow}b_{\pi},c_{\pi},b''_{\pi}$ sont équivalentes, car les $i'\backslash I$ sont \unsure{ordonnées|ordonnés}, et il est clair qu'\uline{alors} $b''_{\pi}\unsure{\Longleftrightarrow|\Longrightarrow}\unsure{b^{\mathrm{s}}_{\pi}|b_{\pi}}$\unsure{.}

% l. 16
\deleted{\uline{4.9.7} Démonstration de 4.9.\unsure{3|2}}
% archive sheet number: 325
\scanpage{327}\mspage{243}
% l. 1–2
4.10. \uline{Le théorème de représentation indicielle} : \uline{démonstration}.

% l. 3–5
On va d'abord reformuler le th. 4.9.3. sous forme plus générale, en termes de catégories \uline{cofibrées}. Soit
% l. 6
\[
 \mathscr D\,\deleted{\unsure{D_j}}\longrightarrow J \tag*{(4.10.1)}
\]
% l. 7–11
un foncteur cofibrant, où $J$ est équivalente à une petite catégorie, et de $\overline{\mathrm m}$ pour les catégories fibres \deleted{$\unsure{C_j|D_j}$} $\mathscr D_j$ $(j\in\operatorname{Ob}J)$. On \deleted{suppose qu'on a} \inserted{se donne} un cardinal \inserted{infini} $\unsure{\pi|\kappa}$ \deleted{tel que}
\[
 \deleted{\operatorname{Card}\unsure{\mathrm{réd}|\mathrm{Fl}\,\mathrm{de}|\mathrm{Flct}}\,J\leq\pi}. \tag*{\deleted{(4.10.2)}}
\]
% l. 12–13
Soit
\[
 \unsure{M|\mathscr M}_j\deleted{\unsure{j}}=\operatorname{Ind}_{\pi}(\mathscr D_j)
 \qquad(j\in\operatorname{Ob}J), \tag*{(4.10.2)}
\]
% l. 14–18
alors les $M_j$ sont les fibres d'une catégorie cofibrée \inserted{$\unsure{\mathscr N|\mathscr M}$} sur $J$, que je désigne aussi par $\operatorname{Ind}_{\pi}(\mathscr D/J)$, donc on pose
\[
 \mathscr N=\operatorname{Ind}_{\pi}(\mathscr D/J). \tag*{(\unsure{4.10.3})}
\]
% l. 19–23
On \deleted{\unsure{peut}} a donc une inclusion \inserted{pl. fid.} de catégories cofibrées sur $J$
\[
 \mathscr D\hookrightarrow\mathscr N,
\]
d'où une \unsure{inclusion|induction}
\[
 \underline{\Gamma}(\mathscr D/J)\longrightarrow\underline{\Gamma}(\mathscr N/J). \tag*{(4.10.4)}
\]
% archive sheet number: 326
\scanpage{328}\mspage{244}
% l. 1–5
Dans le cas où $\mathscr D$ est de la forme $C\times J$, alors $\unsure{\mathscr N|\mathscr M}$ est de la \unsure{formé|forme} \deleted{$\unsure{\operatorname{Ind}}$} $M\times J$, avec $M=\operatorname{Ind}_{\unsure{\pi|\kappa}}(C)$, et l'inclusion (4.10.3) n'est autre \deleted{que} (à isom près) que
\[
 \underline{\mathrm{Hom}}(J,C)\longrightarrow\underline{\mathrm{Hom}}(J,M).
\]

% l. 6–11
\unsure{Revenant} au cas général, on voit encore que $1^\circ)$ Le deuxième \unsure{mb|mbr} de (4.10.3) est stable par $\underline{\lim}$ grandes devant $\pi$, lesquelles se calculent fibre \unsure{à|:} fibre, et $2^\circ)$ Le premier membre de (4.10.3) est formé d'objets $\pi$-accessibles du second.

% l. 12–14
D'où un foncteur \uline{pleinement} fidèle
\[
 \operatorname{Ind}_{\pi}\bigl(\underline{\Gamma}(\mathscr D/J)\bigr)\hookrightarrow
 \underline{\Gamma}(\mathscr N/J). \tag*{(4.10.4)}
\]

% l. 15–19
\uline{Théorème} 4.10.1. Supposons qu'on ait
\[
 \operatorname{Card}\unsure{\mathrm{red}|\mathrm{réd}}\,J\leq\pi, \tag*{(4.10.1.1)}
\]
et qu'il existe un ordinal limite $\alpha$, tel que \deleted{1) pour tout ordinal $\beta<\alpha$, on ait} \uline{1)} $\operatorname{Card}\deleted{\beta}\alpha\leq\pi$, et \uline{2)} les limites inductives \unsure{de|des}
% archive sheet number: 327
\scanpage{329}\mspage{245}
% l. 1–6
type $I_{\alpha}=\{\beta\mid\beta<\alpha\}$ \unsure{sont|soient} représentables dans les catégories $\mathscr D_j$ $(j\in\operatorname{Ob}J)$, et pour tte flèche $\alpha:j\unsure{\mathbin{-}|\longrightarrow}j'$ \unsure{de|dans} $J$, $\alpha^*:\mathscr D_{j'}\longrightarrow\mathscr D_j$, commute auxdites limites. Alors le foncteur canonique pl. fid. (4.10.4) est une équivalence de catégories.

% l. 7–8
Ceci va être contenu dans le résultat \inserted{un peu} plus précis :

% l. 9–12
\uline{Corollaire} 4.10.2. Les hypothèses étant celles \unsure{de|des} 4.10.1, \deleted{soit $\unsure{\varphi|\phi}\in\operatorname{Ob}\underline{\Gamma}(\unsure{\mathscr N|\mathscr M}/J)$. \unsure{Consid}} \inserted{avec en plus}
\[
 \operatorname{Card}J\leq\unsure{\pi|\kappa} \tag*{(4.10.2.1)}
\]
(qui \deleted{\unsure{résulte de}} \inserted{équivaut à} (4.10.1.1) si $J$ est réduite, cas auquel \unsure{on|cas} peut évidemment se réduire dans 4.10.1). Soit $\varphi\in\operatorname{Ob}\underline{\Gamma}(\mathscr N/J)$, et considérons le foncteur canonique \uline{fidèle} :
% l. 16–17
\[
 \underline{\Gamma}(\mathscr D/J)_{/\varphi}\longrightarrow
 \prod_{j\in J} C_j/\varphi(j). \tag*{(4.10.2.\unsure{2|7})\unsure{.}}
\]
% l. 18–20
Alors la catégorie source \deleted{\unsure{de}} $\underline{\Gamma}(\mathscr D/J)_{/\varphi}$ est grde devant $\pi$, et le foncteur envisagé est cofinal\unsure{.}
% archive sheet number: 328
\scanpage{330}\mspage{246}
% l. 1–4
\unsure{De ceci|D'ici}, le th. 4.10.1 résulte de façon essentiellement \unsure{verbale}, tout comme dans le cas particulier 4.9.3 précédent (cf page 223). \deleted{\unsure{Pour}}
% slanted
\margin{un peu court ! Il faudrait expliciter l'argument ou en faire \unsure{un}, il y en a pour \unsure{deux} \unsure{ou} trois pages\unsure{!}}

% l. 5–11
Pour établir le corollaire, comme la catégorie-but de (4.10.2.2) est \unsure{grde|gde} devant $\unsure{\pi|\kappa}$, comme produit de telles catégories (cf 4.9.6.6 page 228), \inserted{et le foncteur est fidèle,} \inserted{il suffit} en vertu de 4.9.6.3, $2^\circ$, de vérifier les conditions $\mathrm F1$ et $\mathrm F3_{\pi}$ \deleted{\unsure{pour}} (page 225) pour le foncteur envisagé.

% l. 12–15
On va interpréter ce foncteur (4.10.2.2) en termes d'une catégorie cofibrée
\[
 \mathscr D_{\unsure{\varphi|\phi}}\longrightarrow J, \tag*{(4.10.5)}
\]
dont les fibres sont les $\mathscr D_j/\varphi(j)$.
% l. 16–19
On note que si $\alpha:j\unsure{\mathbin{-}|\longrightarrow}j'$ est une flèche de $J$, alors on a
\[
 \left\{\begin{aligned}
 \alpha_* &: C_j\unsure{\mathbin{-}|\longrightarrow} C_{\unsure{j'|j}},\\
 &\text{et }\alpha_*(\varphi(j))\longrightarrow\varphi(j')
 \end{aligned}\right.
\]
% archive sheet number: 329
\scanpage{331}\mspage{247}
% l. 1–3
d'où un foncteur composé
\[
 C_j/\unsure{\varphi|\phi}(j)\unsure{\mathbin{-}|\longrightarrow}
 C_{j'}/\alpha^{\unsure{*|\times}}(\varphi(j))\unsure{\mathbin{-}|\longrightarrow}
 C_{j'}/\varphi(j'),
\]
avec transitivité \unsure{à|:} \unsure{Iso.|iso.} can. près etc.
% l. 4–8
On note que l'on a un isom. can.
\[
 \underline{\Gamma}(\mathscr D_{\varphi}/\deleted{\unsure{\mathscr D|J}}\unsure{J|\mathbb J})\simeq
 \underline{\Gamma}(\mathscr D/J)_{/\varphi}
\]
de sorte que le foncteur \inserted{fidèle} canonique (4.10.2.2.) s'identifie au foncteur
\[
 \underline{\Gamma}(\mathscr D_{\varphi}/J)\longrightarrow
 \prod_{j\in\operatorname{Ob}J}(\mathscr D_{\varphi})_j\,.\unsure{.} \tag*{(4.10.6)}
\]

% l. 9–11
Le corollaire 4.10.2 sera alors une conséquence du résultat plus général \deleted{\unsure{s}} suivant :

% l. 12–15
\uline{Corollaire} 4.10.3. Soit $\unsure{\mathcal E|\mathscr E}\longrightarrow J$ un\unsure{e} \deleted{catégorie} \inserted{foncteur} \unsure{cofibrant|cofibrante}, et $\unsure{\pi|\kappa}$ un cardinal infini. On suppose \deleted{$\mathcal E,J$ \unsure{équiv} \inserted{que $\mathcal E,J$ équiv. à de} petites catégories, et de plus.}
% l. 16–20

a) $\operatorname{Card}J\leq\pi$.

b) Les catégories fibres $\mathcal E_j$ sont équivalentes à des petites catégories, et grdes devant $\pi$.

c) \deleted{\unsure{Elles}} \unsure{Il existe} un ordinal \deleted{infini} \inserted{limite} $\alpha$, tel que $1^\circ)$ \deleted{$\beta<\alpha$ \unsure{induit que|implique}} $\operatorname{Card}\deleted{\beta}\alpha\leq\pi$, et
% archive sheet number: 330
\scanpage{332}\mspage{248}
% l. 1–4
$2^\circ)$ Les limites inductives de type $I_{\unsure{\alpha'|\alpha}}$ sont représentables dans les catégories $\unsure{\mathcal E|\mathscr E}_j$, et les foncteurs \unsure{cochgt|cochgt.} de base $\unsure{\alpha|\lambda}_*:\mathcal E_j\longrightarrow\mathcal E_{j'}$ y commutent.
% l. 4–11
[ C'est le cas p. ex. si les limites inductives filtrantes \uline{dénombrables} sont repr. dans les $\mathcal E_j$, et si les $\alpha_*$ y commutent ; ou s'il existe un cardinal $\unsure{\pi|\kappa}_0\leq\pi$, et un cardinal $\pi'_0>\pi_0$, tel que les limites inductives grandes devant $\pi_0$ et de cardinal $\leq\pi'_0$ soient repr. dans les $\mathcal E_j$, et que les $\alpha_*$ y commutent$\unsure{]|)|{}}$ ].

% l. 12–16
\margin{$\unsure{]\!]|]|\parallel}$}\unsure{Sous|Dans} ces conditions, le foncteur canonique fidèle
\[
 \underline{\Gamma}(\mathcal E/J)\longrightarrow\prod_{j\in\operatorname{Ob}J}\mathcal E_j \tag*{(4.10.3.1)}
\]
est cofinal, et la catégorie source est grde devant $\pi$.

% l. 17–19
Il est \inserted{$\pm$} immédiat que les hyp. de 4.10.1 impliquent que celles de 4.10.3 sont satisfaites pour $\mathcal E:\mathscr D_{\unsure{\varphi|\phi}}\longrightarrow J$ (4.10.5), Pour
% archive sheet number: 331
\scanpage{333}\mspage{249}
% l. 1
ceci, on utilise le

% l. 2–5
\uline{Lemme} 4.10.4. \unsure{Soit|Soient} $C$ une catégorie, et $\unsure{\pi|\kappa}$ un cardinal infini. Alors le foncteur d'inclusion $C\longrightarrow\inserted{M\overset{\unsure{\text{def}|\text{déf}}}{=}}\operatorname{Ind}_{\pi} C$ commute aux $\varinjlim$ de cardinal $\leq\pi$.
\margin{\rule{0.4pt}{4\baselineskip}}

% l. 6–10
Soit en effet $I$ telle que $\operatorname{card}I\leq\pi$, et \deleted{\unsure{$i\in I$}} $(x_i)_{i\in I}$ un syst. inductif\unsure{$_{\pi}$} sur $C$, y ayant une $\unsure{\varinjlim|\underline{\lim}}$ $x$. Soit d'autre part $X=\underset{I'}{\varinjlim}\,\unsure{\xi|\zeta}_{i'}$ un objet de $\operatorname{Ind}_{\pi}\unsure{(|C}$, avec $I'$ grde devant $\pi$.

% l. 11–15
On a donc
\[
 \begin{aligned}
 \mathrm{Hom}_{M}(x,X)&\simeq\underset{I'}{\varinjlim}\,\mathrm{Hom}(x,\xi_{i'})
 \simeq\underset{I'}{\varinjlim}\,\underset{I}{\varprojlim}\,\mathrm{Hom}(x_i,\xi_{i'}\unsure{)}\\
 &\overset{\text{SGA4 I 9.8}}{\simeq}\underset{I}{\varprojlim}\,
 \underbrace{\underset{I'}{\varinjlim}\,\mathrm{Hom}(x_i,\xi_{i'})}_{\simeq\,\mathrm{Hom}(x_i,X)}\\
 &\simeq\underset{I}{\varprojlim}\,\mathrm{Hom}(x_i,X),
 \end{aligned}
\]
% l. 16–17
ce qui prouve que $x$ est aussi une $\varinjlim$ des $x_i$ dans $M$ \unsure{tout entier}.
% archive sheet number: 332
\scanpage{334}\mspage{250}
% l. 1–4
Ainsi, on est réduit à prouver 4.10.3. Appliquant 4.9.6. $2^\circ$, on est ramené à prouver, pour le foncteur envisagé, les conditions $\mathrm F1$ et $\mathrm F3_{\unsure{\pi|\kappa}}$ (page 225).

% l. 5–7
Dans les arguments qui suivent, on va utiliser les conditions a), b) de 4.10.3, mais non, pour l'instant, la condition c.

% l. 8–15
\uline{Lemme} 4.10.5. Sous les conditions de 4.10.3 a) b) (à l'exclusion de c), pour que (4.10.3.1) soit cofinal et $\Gamma(\unsure{\mathcal E|\mathscr E}/J)$ grde devant $\pi$, il faut et il suffit que le foncteur envisagé satisfasse à la condition $\mathrm F1$ \deleted{\unsure{et à sa variante}} \inserted{rappelée ci-dessous,} et à la variante $\mathrm F'1$ de celle-ci \unsure{énoncée|donnée} ci-dessous :

% l. 16–19
$\unsure{\bullet|\cdot}\ \mathrm F1)$ Pour \unsure{tout|tt} système $(X(j))_{j\in\operatorname{Ob}J}$ d'éléments $X(j)\in\operatorname{Ob}\mathcal E_j$, existe une \uline{section} $\unsure{G|E}$ de $\mathcal E$ sur $J$, et un système de flèches
\[
 X(j)\longrightarrow G(j)\qquad (j\in\operatorname{Ob}J). \tag*{(4.10.\unsure{5|3}.1)}
\]

% l. 20–21
$\mathrm F'1)$ Sous les conditions de $\mathrm F1$ \unsure{précédent|précédentes}, supposons de plus une section $F$ de $\mathcal E$ sur $J$,
% archive sheet number: 333
\scanpage{335}\mspage{251}
% l. 1–3
\deleted{\unsure{et}} et un système de flèches
\[
 F(j)\unsure{\mathbin{-}|\longrightarrow} X(j)\qquad (j\in\operatorname{Ob}J).
\]
\deleted{\unsure{Alors}} \inserted{alors, $G$ et} le système de flèches (4.10.5.1) peut être choisi de telle façon que les flèches composées
% l. 4–8
\[
 F(j)\unsure{\mathbin{-}|\longrightarrow} X(j)\unsure{\mathbin{-}|\longrightarrow}G(j)
\]
définissent un hom de sections
\[
 F\longrightarrow G
\]
\unsure{(de la} catégorie $\underline{\Gamma}(\unsure{\mathcal E|\mathscr E}/J)$).

% slanted
\margin{\uline{NB} Les conditions $\mathrm F1$ et $\unsure{\mathrm F'1|\mathrm F1}$ équivalent à dire \deleted{que}\inserted{que} le foncteur envisagé satisfait les conditions $\mathrm F1$ et $\mathrm F3$, ou \unsure{encore|mieux} \deleted{que} (en vertu de \unsure{4.9.6.3}, $1^\circ$) que $\underline{\Gamma}(\deleted{\unsure{\mathscr D|\mathcal E}}\mathcal E/J)$ est \uline{filtrante} (sans plus), et le foncteur \inserted{fidèle} (4.10.3.1) cofinal : Cela implique déjà que $\underline{\Gamma}(\mathcal E/J)$ est grde devant $\unsure{\pi|\kappa}$\unsure{.}}

% l. 9–12
\uline{Dém} Il faut voir que $\mathrm F1\mathbin{\&}\mathrm F'1$ équivaut à $\mathrm F1\mathbin{\&}\mathrm F3_{\pi}$. L'implication $\Longleftarrow$ est \inserted{$\pm$} immédiate, c'est de $\Longrightarrow$ que nous avons besoin, donc
\[
 \mathrm F1\mathbin{\&}\mathrm F'1\Longrightarrow\mathrm F3_{\pi}.
\]

% l. 13–19
Pour ceci, \deleted{rappelons} réénonçons $\mathrm F3_{\pi}$ dans le contexte présent. On se donne une famille $(F_{\alpha})_{\alpha\in A}$, $\operatorname{card}A\leq\pi$, d'objets de $\underline{\Gamma}(\mathcal E/J)$, un objet $F$ de la $\overline{\mathrm m}$ \unsure{catégor.|catégorie}, et \deleted{pour} tout $j$ dans $\operatorname{Ob}J$, un système de flèches
\[
 F_{\alpha}(j)\xrightarrow{u_{\alpha}(j)} F(j),
\]
mais sans supposer que pour $\alpha$ fixé, $j$ variable,
% archive sheet number: 334
\scanpage{336}\mspage{252}
% l. 1–4
cela définisse un hom. de sections $F_{\alpha}\unsure{\mathbin{-}|\longrightarrow}F$. On cherche une section $G$ de $\unsure{\mathcal E|\mathscr E}$ sur $J$, et un hom de sections $F\longrightarrow G$, tel\deleted{\unsure{le}} que pour tt $\alpha$, le système \unsure{des|de} \unsure{composés|composées}
% l. 5–6
\[
 F_{\alpha}(j)\unsure{\mathbin{-}|\longrightarrow} F(j)\longrightarrow G(j)
\]
\deleted{\unsure{provienne de}} définisse un hom de sections $F_{\alpha}\longrightarrow G$.

% l. 7–11
Soit donc $\unsure{\lambda|\alpha}:j\longrightarrow j'$ une flèche de $J$, considérons les carrés \unsure{(}\uline{non commutatifs} !\unsure{,}
\diagram{\begin{tikzcd}[ampersand replacement=\&,column sep=small,row sep=small]
 \& F_{\alpha}(j)\arrow[rr,"u_{\alpha}(j)"]\arrow[dd,"F_{\alpha}(\lambda)"'] \& \& F(j)\arrow[dd,"F(\lambda)"]\\
 (\alpha)\& \& ! \&\\
 \& F_{\alpha}(j')\arrow[rr,"u_{\alpha}(j')"]\& \& F(j')
\end{tikzcd}}
(indexés par $\alpha\in A$). Pour $j'$ fixé dans $J$, le nombre de ces carrés est majoré par $\operatorname{Card}\operatorname{Ob}J_{/j'}\times\operatorname{Card}A\leq\unsure{\pi|\kappa}\mathbin{.}\pi=\pi$.

% l. 12–17
Il s'ensuit, comme $\mathcal E_{j'}$ est grande devant $\pi$, qu'il existe un objet $X(j')$ de $\mathcal E_{j'}$, et une flèche \inserted{$v(j'):$} $F(j')\unsure{\mathbin{-}|\longrightarrow} X(j')$ de telle façon que pour tout $\alpha$\inserted{$\in A$}, et tout $\lambda\in\operatorname{Ob}J_{/j'}$, le carré \deleted{\unsure{composé de la}} \inserted{\unsure{déduit|ci-dessus} de \unsure{ce|ces}}
% archive sheet number: 335
\scanpage{337}\mspage{253}
% l. 1–4
$(\alpha)$ en composant avec $F(j')\unsure{\mathbin{-}|\longrightarrow}X(j')$ soit commutatif. On \deleted{choisit} fait ce choix pour tt objet $j'$ de $J$. Appliquons alors $\mathrm F'1$ \inserted{à $F$ et} au système de flèches
% l. 5–7
$F(j)\xrightarrow{v(j)}X(j)$, on trouve $G$ et un système de flèches \deleted{$\unsure{w(j)|v(j)}$}
\[
 w(j):X(j)\unsure{\mathbin{-}|\longrightarrow}G(j)
\]
% l. 8–11
de telle façon que \deleted{pour tte flèche} les flèches composées
\diagram{\begin{tikzcd}[ampersand replacement=\&,column sep=small,row sep=small]
 w(j)v(j):\&F(j)\arrow[rr]\arrow[dr,dashed]\& \&G(j)\\
 \& \&X(j)\arrow[ur,dashed]\&
\end{tikzcd}}
définissent un hom de sections\unsure{,|{}}. Je dis que cet hom. \unsure{à|a} la propriété voulue pour $\mathrm F3_{\unsure{\pi|\kappa}}$, \unsure{i.e|i.e.|vz} que pour tt $\alpha\in A$, le système des composés
% l. 12–16
\[
 F_{\alpha}(j)\unsure{\mathbin{-}|\longrightarrow} F(j)\unsure{\mathbin{-}|\longrightarrow}G(j)
\]
(\deleted{\unsure{pour}} $j\in\operatorname{Ob}J$ variable) définit un hom de sections $F_{\alpha}\longrightarrow G$. En effet, \deleted{\unsure{évident}} \inserted{il suffit} pour cela de \unsure{considérer}, pour une flèche $\lambda$ de $J$, le
% archive sheet number: 336
\scanpage{338}\mspage{254}
% l. 1–4
diagramme
\diagram{\begin{tikzcd}[ampersand replacement=\&,column sep=normal]
 F_{\alpha}(j)\arrow[r,"u_{\alpha}(j)"]\arrow[d,"F_{\alpha}(\lambda)"',no head]\arrow[dr,phantom,"!" description]
 \& F(j)\arrow[r,"v(j)",no head]\arrow[d,"F(\lambda)"']
 \& X(j)\arrow[r,"w(j)"]\arrow[d,phantom,"\text{comm.}" description]
 \& G(j)\arrow[d,"G(\lambda)"]\\
 F_{\alpha}(j')\arrow[r,"u_{\alpha}(j')",no head]\& F(j')\arrow[r,"v(j')",no head]\&
 X(j')\arrow[r,"w(j')",no head]\&G(j')\,.
\end{tikzcd}}

% l. 5–8
Le lemme 4.10.5. étant prouvé, \unsure{tout|il} revient donc à prouver la condition $\mathrm F1$, et son complément $\mathrm F'1$, qui se \unsure{prouve|prouvent} dans la $\overline{\mathrm m}$ haleine.

% l. 9–10
4.10.6. On se donne donc le système d'objets
\[
 X(j)\in\operatorname{Ob}\unsure{\mathcal E|\mathscr E}_j\qquad (j\in\operatorname{Ob}J), \tag*{(4.10.6.1)}
\]
% l. 11–13
et le cas échéant, un objet $F$ de $\underline{\Gamma}(\mathcal E/J)$ et un système de flèches
\[
 F(j)\overset{u(j)}{\unsure{\mathbin{-}|\longrightarrow}}X(j), \tag*{(4.10.6.2)}
\]
% l. 14–16
et on se propose de construire $G$ dans $\underline{\Gamma}(\mathcal E/J)$, et un système de flèches
\[
 X(\unsure{j|j'})\xrightarrow{v(j)}G(j), \tag*{(4.\unsure{10.6|10.26}.3)}
\]
% l. 17–19
de telle façon que (si on a la donnée supplémentaire des $u(j)$) les flèches composées $\deleted{\unsure{X}}F(j)\overset{u(j)}{\unsure{\mathbin{-}|\longrightarrow}}X(j)\overset{v(j)}{\unsure{\mathbin{-}|\longrightarrow}}G(j)$
% archive sheet number: 337
\scanpage{339}\mspage{255}
% l. 1
soient fonctorielles en $j$.

% l. 2–8
On va faire une construction transfinie dans les ordinaux $\alpha$ tels que
\[
 \operatorname{Card}\alpha\leq\unsure{\pi|\kappa}
\]
d'un système inductif
\[
 (X_{\alpha}(\unsure{j|g}))_{\alpha\deleted{\unsure{j|J}}}
\]
d'objets de $\prod_{j\in\operatorname{Ob}J}\unsure{\mathcal E|\mathscr E}_j$. On va d'abord ignorer la donnée éventuelle (4.10.6.2).

% l. 9–12
On posera
\[
 \deleted{X_0(j)=X_j}\qquad X_0(j)=X(j), \tag*{(4.10.6.4)}
\]
et on va donner la prescription pour passer de $X_{\alpha}$ à $X_{\alpha+1}$. La construction \deleted{\unsure{sera telle}} \unsure{inclura} aussi, pour tte flèche $\unsure{\lambda|\alpha}$ \inserted{$:j'\unsure{\mathbin{-}|\longrightarrow}j\deleted{\unsure{'}}$} de $J$ et tout $\alpha$, la donnée de
% l. 13–16
\[
 X_{\alpha}(j')\xrightarrow{\unsure{t_{\alpha}|\tau_{\alpha}}(\lambda)}X_{\alpha+1}(j\deleted{\unsure{'}}) \tag*{(4.10.6.5)}
\]
\margin{flèche (dans $\mathcal E$) \unsure{au} dessus de $\lambda$}
et ces flèches, et les flèches de transition $X_{\alpha}(j)\overset{\unsure{\tau|T}_{\alpha}(j)}{\unsure{\mathbin{-}|\longrightarrow}} X_{\deleted{\unsure{j}}\unsure{\alpha+1|\alpha,}}(j)$, doivent
% archive sheet number: 338
\scanpage{340}\mspage{256}
% l. 1–3; brace in left margin joins (i), (ii)
satisfaire aux \deleted{trois} conditions suivantes :

\noindent\llap{\smash{\raisebox{-0.8\baselineskip}{$\Biggl\{$}}\quad}(i) Si $\lambda$ est une flèche identique, \inserted{$\mathrm{id}_j:j\unsure{\mathbin{-}|\longrightarrow}j$} alors $\boxed{\deleted{\unsure{t}}t_\alpha(\lambda)=\unsure{\tau|T}_\alpha(j)}$

(ii) Dans le diagramme
% l. 4–6
\diagram{\begin{tikzcd}[ampersand replacement=\&,column sep=small]
 X_\alpha(j') \arrow[r,no head,"{\tau_\alpha(j')}"] \arrow[dr,"{t_\alpha(\lambda)}"] \arrow[d,phantom,"{\deleted{\downarrow}}" description]
 \& X_{\alpha+1}(j') \arrow[r,no head] \arrow[dr,"{t_{\alpha+1}(\lambda)}"] \arrow[d,phantom,"{\bullet}" description]
 \& X_{\alpha+2}(j') \arrow[r,no head] \arrow[dr] \& X_{\alpha+3}(j') \arrow[r] \arrow[dr,dotted,no head] \& {}\\
 X_\alpha(j\deleted{\unsure{'}}) \arrow[r,no head] \& X_{\alpha+1}(j\deleted{\unsure{'}}) \arrow[r,no head,"{\tau_{\alpha+1}(j\deleted{\unsure{'}})}"]
 \& X_{\alpha+2}(j\deleted{\unsure{'}}) \arrow[r,very thick,"{\tau_{\alpha+2}(j\deleted{\unsure{'}})}"]
 \& X_{\alpha+3}(j\deleted{\unsure{'}}) \arrow[r] \& {}
\end{tikzcd}}
% l. 7–9
$\tau_{\alpha+2}(j\deleted{\unsure{'}})$ ``rend le carré $\bullet$ commutatif'', \deleted{$\unsure{\tau_{\alpha+2}(j')}$} i.e. égalise $t_{\alpha+1}(\lambda)\tau_\alpha(j')$ et $\tau_{\alpha+1}(j\deleted{\unsure{'}})t_\alpha(\lambda)$ :
\[
 \boxed{\tau_{\alpha+2}(j)t_{\alpha+1}(\lambda)\tau_\alpha(j')=\tau_{\alpha+2}(j)\tau_{\alpha+1}(j\deleted{\unsure{'}})t_\alpha(\lambda)}
\]
% l. 10–13
(iii) Si on a deux flèches consécutives $\lambda:j''\unsure{\mathbin{-}|\longrightarrow}j'$, $\mu:j'\unsure{\mathbin{-}|\longrightarrow}j\deleted{\unsure{'}}$, on a dans le diagramme
\diagram{\begin{tikzcd}[ampersand replacement=\&,column sep=small]
 X_\alpha(j'') \arrow[r] \arrow[dr,"{t_\alpha(\lambda)}"] \arrow[ddr,"{t_\alpha(\mu\lambda)}" near end]
 \& X_{\alpha+1}(j'') \arrow[r,no head] \arrow[dr] \& X_{\alpha+2}(j'') \arrow[r,no head] \arrow[dr] \& X_{\alpha+3}(j'')\\
 X_\alpha(\deleted{\unsure{j'}}) \arrow[r,phantom,"{\deleted{\longrightarrow}}" description] \arrow[dr,phantom,"{\deleted{\searrow}}" description]
 \& X_{\alpha+1}(j') \arrow[r,very thick,no head] \arrow[dr,"{t_{\alpha+1}(\mu)}"]
 \& X_{\alpha+2}(j') \arrow[r,no head] \arrow[dr] \& X_{\alpha+3}(j')\\
 X_\alpha(j\deleted{\unsure{'}}) \arrow[r] \& X_{\alpha+1}(j\deleted{\unsure{'}}) \arrow[r,no head,"{\tau_{\alpha+1}(j\deleted{\unsure{'}})}"]
 \& X_{\alpha+2}(j\deleted{\unsure{'}}) \arrow[r,"{\tau_{\alpha+2}(j\deleted{\unsure{'}})}"] \& X_{\alpha+3}(j\deleted{\unsure{'}})
\end{tikzcd}}
% l. 14
\[
 \boxed{\tau_{\alpha+2}(j\deleted{\unsure{'}})t_{\alpha+1}(\mu)t_\alpha(\lambda)=\tau_{\alpha+2}(j\deleted{\unsure{'}})\deleted{\unsure{t}}\unsure{\tau|t}_{\alpha+1}(j)t_\alpha(\mu\lambda)}
\]
% archive sheet number: 339
